\documentclass{article}
\usepackage{/globals/laws}
\usepackage{/globals/anki}
\ankiprefix{gesetzgebungskompetenzen}
\ankitopic{Gesetzgebungskompetenzen}
\ankishow
\begin{document}
\section{Gesetzgebungskompetenzen} 
Ein Staat hat das Recht Gesetze in Bezug auf einen Bereich, die Gesetzgebungsmaterie, zu erlassen, verändern oder aufzuheben.

In Deutschland haben sowohl der Bund als auch die Länder eine Gesetzgebungskompetenzen, in unterschiedlichen Bereichen. Wird ein Gesetz auf der falschen Eben erlassen, so ist es formal verfassungswidrig.

Dies wird in \arts{70}{74} \GG behandelt.
\ankinote{gg}{GG}{Art. 70-74}

\begin{description}
 \item{Ausschließliche Gesetzgebungskompetenz} Eigentlich ist nur der Bund zuständig, der Bund kann aber in Gesetzen Länder dazu ermächtigen eigene Gesetze zu erlassen.
 \ankinote{ausschließlich}{Ausschließlich}{Nur Bund, kann Länder ermächtigen}
 \item{Konkurrierende Gesetzgebungskompetenz} Der Bund hat vorrang und darf Gesetze erlassen, Länder dürfen aber auch Gesetze erlassen solange der Bund noch keine Gesetze erlassen hat, es hat nicht abschließend ist, oder es ein Abweichungsbefugnis gibt.
 \begin{description}
  \item{Bedingungskompetenz} Der Bund darf nur ein Gesetz erlassen solange bestimmte Bedingungen erfüllt sind.
  \item{bedingungslose Kompetenz} Der Bund darf bedingungslos ein Gesetz erlassen.
  \item{Abweichungskompetenz} Länder dürfen in bestimmten Bereichen Gesetze erlassen die vom Bundesrecht abweichen.
 \end{description}
 \ankinote{konkurrierend}{Konkurrierend}{Bund darf, Länder wenn Bund nicht abschließend. Bedingungs-, bedingungslose-, Abweichungs-}
 \ankinote{bedingung}{Bedingungskompetenz}{Wenn bestimmte Bedingungen erfüllt sind}
 \ankinote{bedingungslos}{Bedingungslos}{Ohne Bedingungen}
 \ankinote{abweichung}{Abweichungskompetenz}{Länder in Bereichen von Bund abweichen}
 \item{Ungeschriebene Kompetenzen} 
 \begin{description}
  \item[Natur der Sache] Manche Gesetze ergeben nur auf einer Ebene Sinn.
  \item[Annexkompetenz] Eine Sache für die keine Kompetenz besteht, die aber so ähnlich ist wie eine Sache für die Kompetenz besteht, muss auch geregelt werden.
 \end{description}
 \ankinote{ungeschrieben}{Ungeschriebene Kompetenzen}{Natur der Sache, Annexkompetenz}
 \ankinote{natur}{Natur der Sache}{Gibt nur auf der Ebene Sinn}
 \ankinote{annex}{Annexkompetenz}{Ist so ähnlich wie etwas mit Kompetenz,Sachzusammenhangs-}
\end{description}


\end{document}